\subsection{Operational framework}

Building on these distinctions, this dissertation adopts an operational framework for describing and comparing different implementations of Climate Services. The framework identifies the elements and relationships through which climate information may become relevant to decisions without defining how a Climate Service should be designed.

These relationships cannot be represented adequately as a linear pipeline from a single producer to a final user. User requirements may affect which data are collected and how they are processed, while feedback may lead to revisions of products, methods, or the initial formulation of the problem. Moreover, the same organization may produce information in one context, provide it in another, and use information generated elsewhere. Producers, providers, intermediaries, decision-makers, and other stakeholders consequently form a continuous and context-dependent network of overlapping roles \citep{weichselgartner_evolving_2019}.

The framework represents this relational structure as a \textit{value web}. Although the climate service landscape may be described in terms of a value chain, its activities commonly involve iterations, feedback loops, and connections that a linear sequence cannot capture \citep{hewitt_climate_2021}. A value-web perspective accommodates these interactions without assigning actors permanently to a single position. It also situates them within an enabling environment shaped by governance and regulatory arrangements, funding mechanisms, technical infrastructures, institutional capacities, and other contextual conditions.

Within this value web, the framework distinguishes four analytical elements: \textit{inputs}, including climate and non-climate data, expertise, resources, and user requirements; \textit{processes}, including analysis, translation, tailoring, interaction, and, where appropriate, co-production; \textit{outputs}, primarily the Climate Products through which data and information are communicated; and \textit{outcomes}, concerning whether and how these outputs are used and influence decisions, actions, or decision-making capacities. This structure is consistent with the Climateurope2 framework for assessing Climate Services across input, co-production, output, and outcome \citep{villwock_literature_2023}. The four elements are not consecutive stages, but interconnected entry points for analysing the activities and relationships that constitute a Climate Service.